% Options for packages loaded elsewhere
\PassOptionsToPackage{unicode}{hyperref}
\PassOptionsToPackage{hyphens}{url}
%
\documentclass[
]{article}
\usepackage{amsmath,amssymb}
\usepackage{iftex}
\ifPDFTeX
  \usepackage[T1]{fontenc}
  \usepackage[utf8]{inputenc}
  \usepackage{textcomp} % provide euro and other symbols
\else % if luatex or xetex
  \usepackage{unicode-math} % this also loads fontspec
  \defaultfontfeatures{Scale=MatchLowercase}
  \defaultfontfeatures[\rmfamily]{Ligatures=TeX,Scale=1}
\fi
\usepackage{lmodern}
\ifPDFTeX\else
  % xetex/luatex font selection
\fi
% Use upquote if available, for straight quotes in verbatim environments
\IfFileExists{upquote.sty}{\usepackage{upquote}}{}
\IfFileExists{microtype.sty}{% use microtype if available
  \usepackage[]{microtype}
  \UseMicrotypeSet[protrusion]{basicmath} % disable protrusion for tt fonts
}{}
\makeatletter
\@ifundefined{KOMAClassName}{% if non-KOMA class
  \IfFileExists{parskip.sty}{%
    \usepackage{parskip}
  }{% else
    \setlength{\parindent}{0pt}
    \setlength{\parskip}{6pt plus 2pt minus 1pt}}
}{% if KOMA class
  \KOMAoptions{parskip=half}}
\makeatother
\usepackage{xcolor}
\usepackage[margin=1in]{geometry}
\usepackage{color}
\usepackage{fancyvrb}
\newcommand{\VerbBar}{|}
\newcommand{\VERB}{\Verb[commandchars=\\\{\}]}
\DefineVerbatimEnvironment{Highlighting}{Verbatim}{commandchars=\\\{\}}
% Add ',fontsize=\small' for more characters per line
\usepackage{framed}
\definecolor{shadecolor}{RGB}{248,248,248}
\newenvironment{Shaded}{\begin{snugshade}}{\end{snugshade}}
\newcommand{\AlertTok}[1]{\textcolor[rgb]{0.94,0.16,0.16}{#1}}
\newcommand{\AnnotationTok}[1]{\textcolor[rgb]{0.56,0.35,0.01}{\textbf{\textit{#1}}}}
\newcommand{\AttributeTok}[1]{\textcolor[rgb]{0.13,0.29,0.53}{#1}}
\newcommand{\BaseNTok}[1]{\textcolor[rgb]{0.00,0.00,0.81}{#1}}
\newcommand{\BuiltInTok}[1]{#1}
\newcommand{\CharTok}[1]{\textcolor[rgb]{0.31,0.60,0.02}{#1}}
\newcommand{\CommentTok}[1]{\textcolor[rgb]{0.56,0.35,0.01}{\textit{#1}}}
\newcommand{\CommentVarTok}[1]{\textcolor[rgb]{0.56,0.35,0.01}{\textbf{\textit{#1}}}}
\newcommand{\ConstantTok}[1]{\textcolor[rgb]{0.56,0.35,0.01}{#1}}
\newcommand{\ControlFlowTok}[1]{\textcolor[rgb]{0.13,0.29,0.53}{\textbf{#1}}}
\newcommand{\DataTypeTok}[1]{\textcolor[rgb]{0.13,0.29,0.53}{#1}}
\newcommand{\DecValTok}[1]{\textcolor[rgb]{0.00,0.00,0.81}{#1}}
\newcommand{\DocumentationTok}[1]{\textcolor[rgb]{0.56,0.35,0.01}{\textbf{\textit{#1}}}}
\newcommand{\ErrorTok}[1]{\textcolor[rgb]{0.64,0.00,0.00}{\textbf{#1}}}
\newcommand{\ExtensionTok}[1]{#1}
\newcommand{\FloatTok}[1]{\textcolor[rgb]{0.00,0.00,0.81}{#1}}
\newcommand{\FunctionTok}[1]{\textcolor[rgb]{0.13,0.29,0.53}{\textbf{#1}}}
\newcommand{\ImportTok}[1]{#1}
\newcommand{\InformationTok}[1]{\textcolor[rgb]{0.56,0.35,0.01}{\textbf{\textit{#1}}}}
\newcommand{\KeywordTok}[1]{\textcolor[rgb]{0.13,0.29,0.53}{\textbf{#1}}}
\newcommand{\NormalTok}[1]{#1}
\newcommand{\OperatorTok}[1]{\textcolor[rgb]{0.81,0.36,0.00}{\textbf{#1}}}
\newcommand{\OtherTok}[1]{\textcolor[rgb]{0.56,0.35,0.01}{#1}}
\newcommand{\PreprocessorTok}[1]{\textcolor[rgb]{0.56,0.35,0.01}{\textit{#1}}}
\newcommand{\RegionMarkerTok}[1]{#1}
\newcommand{\SpecialCharTok}[1]{\textcolor[rgb]{0.81,0.36,0.00}{\textbf{#1}}}
\newcommand{\SpecialStringTok}[1]{\textcolor[rgb]{0.31,0.60,0.02}{#1}}
\newcommand{\StringTok}[1]{\textcolor[rgb]{0.31,0.60,0.02}{#1}}
\newcommand{\VariableTok}[1]{\textcolor[rgb]{0.00,0.00,0.00}{#1}}
\newcommand{\VerbatimStringTok}[1]{\textcolor[rgb]{0.31,0.60,0.02}{#1}}
\newcommand{\WarningTok}[1]{\textcolor[rgb]{0.56,0.35,0.01}{\textbf{\textit{#1}}}}
\usepackage{longtable,booktabs,array}
\usepackage{calc} % for calculating minipage widths
% Correct order of tables after \paragraph or \subparagraph
\usepackage{etoolbox}
\makeatletter
\patchcmd\longtable{\par}{\if@noskipsec\mbox{}\fi\par}{}{}
\makeatother
% Allow footnotes in longtable head/foot
\IfFileExists{footnotehyper.sty}{\usepackage{footnotehyper}}{\usepackage{footnote}}
\makesavenoteenv{longtable}
\usepackage{graphicx}
\makeatletter
\def\maxwidth{\ifdim\Gin@nat@width>\linewidth\linewidth\else\Gin@nat@width\fi}
\def\maxheight{\ifdim\Gin@nat@height>\textheight\textheight\else\Gin@nat@height\fi}
\makeatother
% Scale images if necessary, so that they will not overflow the page
% margins by default, and it is still possible to overwrite the defaults
% using explicit options in \includegraphics[width, height, ...]{}
\setkeys{Gin}{width=\maxwidth,height=\maxheight,keepaspectratio}
% Set default figure placement to htbp
\makeatletter
\def\fps@figure{htbp}
\makeatother
\setlength{\emergencystretch}{3em} % prevent overfull lines
\providecommand{\tightlist}{%
  \setlength{\itemsep}{0pt}\setlength{\parskip}{0pt}}
\setcounter{secnumdepth}{-\maxdimen} % remove section numbering
\ifLuaTeX
  \usepackage{selnolig}  % disable illegal ligatures
\fi
\usepackage{bookmark}
\IfFileExists{xurl.sty}{\usepackage{xurl}}{} % add URL line breaks if available
\urlstyle{same}
\hypersetup{
  pdftitle={How To Build a Many-Model Forecasting Solution on Snowflake},
  hidelinks,
  pdfcreator={LaTeX via pandoc}}

\title{How To Build a Many-Model Forecasting Solution on Snowflake}
\author{}
\date{\vspace{-2.5em}}

\begin{document}
\maketitle

This guide walks through the complete \textbf{many-model pattern} on
Snowflake: fit hundreds or thousands of independent models in parallel
using Snowpark Container Services (SPCS), persist them to a stage,
register them in Model Registry as a single model version, and run
server-side batch inference via \texttt{run\_batch}.

The accompanying \textbf{reference implementation} lives in
\texttt{internal/doSnowflake/demos/parallel\_forecast/}. The code
snippets below are extracted from that implementation. Lines marked
\texttt{\#\#\ ADAPT} are the parts you replace for your own data and
models.

\subsection{1. Architecture Overview}\label{architecture-overview}

The many-model workflow has four phases:

\begin{verbatim}
 ┌──────────────┐   foreach / SPCS     ┌──────────────┐
 │ SERIES_EVENTS├──────────────────────► .rds models  │
 │  (data table)│  N parallel workers  │  on @STAGE   │
 └──────────────┘                      └──────┬───────┘
                                              │
                                   sfr_log_many_model()
                                              │
                                       ┌──────▼───────┐
                                       │Model Registry│
                                       │ (aggregator) │
                                       └──────┬───────┘
                                              │
                                    sfr_run_batch()
                                              │
                                      ┌───────▼────────┐
                                      │FORECAST_RESULTS│
                                      │  (output table)│
                                      └────────────────┘
\end{verbatim}

\textbf{Two dispatch modes} are available:

\begin{longtable}[]{@{}
  >{\raggedright\arraybackslash}p{(\columnwidth - 4\tabcolsep) * \real{0.2222}}
  >{\raggedright\arraybackslash}p{(\columnwidth - 4\tabcolsep) * \real{0.4074}}
  >{\raggedright\arraybackslash}p{(\columnwidth - 4\tabcolsep) * \real{0.3704}}@{}}
\toprule\noalign{}
\begin{minipage}[b]{\linewidth}\raggedright
Mode
\end{minipage} & \begin{minipage}[b]{\linewidth}\raggedright
Mechanism
\end{minipage} & \begin{minipage}[b]{\linewidth}\raggedright
Best for
\end{minipage} \\
\midrule\noalign{}
\endhead
\bottomrule\noalign{}
\endlastfoot
\textbf{Tasks} & Snowflake Tasks + SPCS jobs; fixed batch of chunks
dispatched upfront & Known workload; all chunks must complete \\
\textbf{Queue} & Hybrid Table queue; dynamic feeder adds chunks as
workers finish & Time-boxed runs; streaming; scale-out/in \\
\end{longtable}

\textbf{The ``screen then refine'' pattern}: run a fast Tasks pass
(e.g.~stepwise ARIMA) over all SKUs, capture accuracy metrics, then rank
the worst performers and feed them into a time-boxed Queue run that
applies a more expensive model contest (ARIMA + ETS + TBATS with holdout
validation). The time budget is spent where it matters most.

\subsection{2. Prerequisites}\label{prerequisites}

\subsubsection{Snowflake objects}\label{snowflake-objects}

\begin{itemize}
\tightlist
\item
  A \textbf{compute pool} with enough nodes (e.g.~\texttt{CPU\_X64\_L},
  4 nodes).
\item
  An \textbf{image repository} with the \texttt{dosnowflake-worker}
  Docker image pushed.
\item
  A \textbf{warehouse} (XS is fine for orchestration SQL).
\item
  The role used must have \texttt{USAGE} on the compute pool and
  \texttt{CREATE\ SERVICE} / \texttt{CREATE\ TASK} privileges.
\end{itemize}

\subsubsection{Local R environment}\label{local-r-environment}

\begin{Shaded}
\begin{Highlighting}[]
\FunctionTok{install.packages}\NormalTok{(}\FunctionTok{c}\NormalTok{(}\StringTok{"snowflakeR"}\NormalTok{, }\StringTok{"foreach"}\NormalTok{, }\StringTok{"iterators"}\NormalTok{, }\StringTok{"forecast"}\NormalTok{,}
                    \StringTok{"ggplot2"}\NormalTok{, }\StringTok{"yaml"}\NormalTok{))}
\end{Highlighting}
\end{Shaded}

\subsubsection{Authentication}\label{authentication}

Configure \texttt{\textasciitilde{}/.snowflake/connections.toml} with a
named profile using key-pair (JWT) auth:

\begin{Shaded}
\begin{Highlighting}[]
\KeywordTok{[my\_demo\_jwt]}
\DataTypeTok{account}   \OperatorTok{=} \StringTok{"MYORG{-}MYACCOUNT"}
\DataTypeTok{user}      \OperatorTok{=} \StringTok{"MY\_USER"}
\DataTypeTok{authenticator} \OperatorTok{=} \StringTok{"SNOWFLAKE\_JWT"}
\DataTypeTok{private\_key\_path} \OperatorTok{=} \StringTok{"\textasciitilde{}/.snowflake/keys/rsa\_key.p8"}
\DataTypeTok{warehouse} \OperatorTok{=} \StringTok{"MY\_WH"}
\DataTypeTok{database}  \OperatorTok{=} \StringTok{"MY\_DB"}
\DataTypeTok{role}      \OperatorTok{=} \StringTok{"MY\_ROLE"}
\end{Highlighting}
\end{Shaded}

\subsection{3. Config and Bootstrap}\label{config-and-bootstrap}

\subsubsection{The YAML config file}\label{the-yaml-config-file}

All environment-specific values live in a single YAML file. The
reference implementation uses
\texttt{snowflaker\_parallel\_spcs\_benchmark\_ak\_small.yaml}:

\begin{Shaded}
\begin{Highlighting}[]
\CommentTok{\#\# ADAPT: replace every value below with your account\textquotesingle{}s objects}
\FunctionTok{context}\KeywordTok{:}
\AttributeTok{  }\FunctionTok{warehouse}\KeywordTok{:}\AttributeTok{ MY\_WAREHOUSE}
\AttributeTok{  }\FunctionTok{database}\KeywordTok{:}\AttributeTok{  MY\_DATABASE}
\AttributeTok{  }\FunctionTok{schema}\KeywordTok{:}\AttributeTok{    SOURCE\_DATA}

\FunctionTok{parallel\_lab}\KeywordTok{:}
\AttributeTok{  }\FunctionTok{clean\_room}\KeywordTok{:}\AttributeTok{ }\CharTok{false}
\AttributeTok{  }\FunctionTok{database}\KeywordTok{:}\AttributeTok{  MY\_DATABASE}
\AttributeTok{  }\FunctionTok{schemas}\KeywordTok{:}
\AttributeTok{    }\FunctionTok{source\_data}\KeywordTok{:}\AttributeTok{ SOURCE\_DATA}
\AttributeTok{    }\FunctionTok{config}\KeywordTok{:}\AttributeTok{      CONFIG}
\AttributeTok{    }\FunctionTok{models}\KeywordTok{:}\AttributeTok{      MODELS}
\AttributeTok{  }\FunctionTok{warehouse}\KeywordTok{:}\AttributeTok{     MY\_WAREHOUSE}
\AttributeTok{  }\FunctionTok{compute\_pool}\KeywordTok{:}\AttributeTok{  MY\_COMPUTE\_POOL}
\AttributeTok{  }\FunctionTok{image\_uri}\KeywordTok{:}\AttributeTok{     }\StringTok{"myorg{-}myacct.registry.../my\_db/config/images/dosnowflake{-}worker:latest"}
\AttributeTok{  }\FunctionTok{dosnowflake\_stage\_name}\KeywordTok{:}\AttributeTok{ DOSNOWFLAKE\_STAGE}
\AttributeTok{  }\FunctionTok{queue\_table}\KeywordTok{:}\AttributeTok{ DOSNOWFLAKE\_QUEUE}
\AttributeTok{  }\FunctionTok{instance\_family}\KeywordTok{:}\AttributeTok{ CPU\_X64\_L}

\AttributeTok{  }\FunctionTok{demo\_forecast\_n\_skus}\KeywordTok{:}\AttributeTok{ }\DecValTok{2000}
\AttributeTok{  }\FunctionTok{demo\_tasks\_chunks\_per\_job}\KeywordTok{:}\AttributeTok{ }\DecValTok{4}
\AttributeTok{  }\FunctionTok{demo\_queue\_n\_workers}\KeywordTok{:}\AttributeTok{ }\DecValTok{4}
\AttributeTok{  }\FunctionTok{demo\_queue\_chunks\_per\_job}\KeywordTok{:}\AttributeTok{ }\DecValTok{8}
\end{Highlighting}
\end{Shaded}

\subsubsection{Loading config and
connecting}\label{loading-config-and-connecting}

\begin{Shaded}
\begin{Highlighting}[]
\FunctionTok{source}\NormalTok{(}\StringTok{"snowflakeR/inst/notebooks/parallel\_spcs\_workflow.R"}\NormalTok{)}

\NormalTok{cfg  }\OtherTok{\textless{}{-}} \FunctionTok{parallel\_lab\_load\_config}\NormalTok{(}\StringTok{"path/to/my\_config.yaml"}\NormalTok{)}
\NormalTok{conn }\OtherTok{\textless{}{-}}\NormalTok{ snowflakeR}\SpecialCharTok{::}\FunctionTok{sfr\_connect}\NormalTok{(}
  \AttributeTok{name =} \StringTok{"my\_demo\_jwt"}\NormalTok{,}
  \AttributeTok{private\_key\_file =} \StringTok{"\textasciitilde{}/.snowflake/keys/rsa\_key.p8"}
\NormalTok{)}
\end{Highlighting}
\end{Shaded}

\subsubsection{Bootstrap DDL}\label{bootstrap-ddl}

\texttt{parallel\_lab\_setup()} runs bootstrap SQL that creates:

\begin{itemize}
\tightlist
\item
  \textbf{3 schemas}: \texttt{SOURCE\_DATA}, \texttt{CONFIG},
  \texttt{MODELS}
\item
  \textbf{Internal stage} \texttt{DOSNOWFLAKE\_STAGE} with
  \texttt{DIRECTORY\ =\ (ENABLE\ =\ TRUE)}
\item
  \textbf{Hybrid Table} \texttt{DOSNOWFLAKE\_QUEUE} (for queue mode)
\item
  \textbf{Tables}: \texttt{TRAINING\_RESULTS},
  \texttt{TRAINING\_METRICS}, \texttt{FORECAST\_RESULTS}
\item
  \textbf{View}: \texttt{MODEL\_INDEX} (parses \texttt{RELATIVE\_PATH}
  from the stage directory table into \texttt{RUN\_ID} +
  \texttt{PARTITION\_KEY})
\end{itemize}

\begin{Shaded}
\begin{Highlighting}[]
\FunctionTok{parallel\_lab\_setup}\NormalTok{(conn, cfg, }\AttributeTok{create\_series =} \ConstantTok{TRUE}\NormalTok{, }\AttributeTok{n\_units =} \DecValTok{2000}\NormalTok{, }\AttributeTok{n\_days =} \DecValTok{1095}\NormalTok{)}
\end{Highlighting}
\end{Shaded}

Set \texttt{create\_series\ =\ TRUE} for synthetic data (generates 3
years of daily observations per SKU with weekly + annual seasonality,
trend changepoints, and occasional outliers), or \texttt{FALSE} if your
source table already exists.

\textbf{ADAPT}: the \texttt{SERIES\_EVENTS} schema is
\texttt{(UNIT\_ID\ VARCHAR,\ OBS\_DATE\ DATE,\ Y\ FLOAT)}. Replace this
with your own data table's DDL.

\subsection{4. The foreach Expression}\label{the-foreach-expression}

The worker expression is the R code that runs \textbf{inside each SPCS
container} for every unit (SKU/store/sensor/etc.). It receives an
injected \texttt{unit\_data} data.frame and a \texttt{unit\_id}
variable.

\subsubsection{Anatomy of a worker
expression}\label{anatomy-of-a-worker-expression}

\begin{Shaded}
\begin{Highlighting}[]
\DocumentationTok{\#\# This runs inside the SPCS container for each unit\_id.}
\DocumentationTok{\#\# \textasciigrave{}unit\_data\textasciigrave{} is automatically injected by doSnowflake — it contains the rows}
\DocumentationTok{\#\# from SERIES\_EVENTS where UNIT\_ID == unit\_id.}

\FunctionTok{suppressPackageStartupMessages}\NormalTok{(}\FunctionTok{library}\NormalTok{(forecast))}

\DocumentationTok{\#\# ADAPT: convert your data to a ts object with the right frequency}
\DocumentationTok{\#\# frequency = 7 for daily data (weekly cycle); use 12 for monthly, 365.25 for}
\DocumentationTok{\#\# daily with annual cycle, etc.}
\NormalTok{ts\_full }\OtherTok{\textless{}{-}} \FunctionTok{ts}\NormalTok{(unit\_data}\SpecialCharTok{$}\NormalTok{Y, }\AttributeTok{frequency =} \DecValTok{7}\NormalTok{)}

\DocumentationTok{\#\# ADAPT: fit your model here (any R model that produces forecasts)}
\NormalTok{model }\OtherTok{\textless{}{-}} \FunctionTok{auto.arima}\NormalTok{(ts\_full, }\AttributeTok{stepwise =} \ConstantTok{TRUE}\NormalTok{)}
\NormalTok{fc    }\OtherTok{\textless{}{-}} \FunctionTok{forecast}\NormalTok{(model, }\AttributeTok{h =} \DecValTok{30}\DataTypeTok{L}\NormalTok{)}
\NormalTok{acc   }\OtherTok{\textless{}{-}}\NormalTok{ forecast}\SpecialCharTok{::}\FunctionTok{accuracy}\NormalTok{(model)}

\DocumentationTok{\#\# Contract: return a list with these fields}
\FunctionTok{list}\NormalTok{(}
  \AttributeTok{unit\_id   =}\NormalTok{ unit\_id,}
  \AttributeTok{model     =} \FunctionTok{as.character}\NormalTok{(fc}\SpecialCharTok{$}\NormalTok{method),}
  \AttributeTok{forecast  =} \FunctionTok{as.numeric}\NormalTok{(fc}\SpecialCharTok{$}\NormalTok{mean),}
  \AttributeTok{n\_obs     =} \FunctionTok{nrow}\NormalTok{(unit\_data),}
  \AttributeTok{model\_obj =}\NormalTok{ model,           }\CommentTok{\# saved as .rds on stage}
  \AttributeTok{aicc      =}\NormalTok{ model}\SpecialCharTok{$}\NormalTok{aicc,}
  \AttributeTok{rmse      =}\NormalTok{ acc[}\DecValTok{1}\NormalTok{, }\StringTok{"RMSE"}\NormalTok{],}
  \AttributeTok{mape      =}\NormalTok{ holdout\_mape,    }\CommentTok{\# see holdout validation below}
  \AttributeTok{mase      =}\NormalTok{ acc[}\DecValTok{1}\NormalTok{, }\StringTok{"MASE"}\NormalTok{]}
\NormalTok{)}
\end{Highlighting}
\end{Shaded}

\textbf{Key contract}: the returned list must include \texttt{unit\_id}
(the partition key) and \texttt{model\_obj} (the fitted model object ---
serialized as \texttt{.rds} on the stage). Additional fields
(\texttt{forecast}, metrics) are collected in the results list.

\subsubsection{Holdout validation}\label{holdout-validation}

For fair cross-method comparison, compute an \textbf{out-of-sample
holdout MAPE} by reserving the last 30 days (or 1/4 of the series) as a
test set:

\begin{Shaded}
\begin{Highlighting}[]
\NormalTok{n }\OtherTok{\textless{}{-}} \FunctionTok{length}\NormalTok{(ts\_full)}
\NormalTok{holdout\_n }\OtherTok{\textless{}{-}} \FunctionTok{min}\NormalTok{(}\DecValTok{30}\DataTypeTok{L}\NormalTok{, }\FunctionTok{as.integer}\NormalTok{(}\FunctionTok{floor}\NormalTok{(n }\SpecialCharTok{/} \DecValTok{4}\NormalTok{)))}

\NormalTok{holdout\_mape }\OtherTok{\textless{}{-}} \FunctionTok{tryCatch}\NormalTok{(\{}
\NormalTok{  ts\_train }\OtherTok{\textless{}{-}} \FunctionTok{window}\NormalTok{(ts\_full, }\AttributeTok{end =} \FunctionTok{time}\NormalTok{(ts\_full)[n }\SpecialCharTok{{-}}\NormalTok{ holdout\_n])}
\NormalTok{  ts\_test  }\OtherTok{\textless{}{-}} \FunctionTok{window}\NormalTok{(ts\_full, }\AttributeTok{start =} \FunctionTok{time}\NormalTok{(ts\_full)[n }\SpecialCharTok{{-}}\NormalTok{ holdout\_n }\SpecialCharTok{+} \DecValTok{1}\NormalTok{])}
\NormalTok{  m\_ho     }\OtherTok{\textless{}{-}} \FunctionTok{auto.arima}\NormalTok{(ts\_train, }\AttributeTok{stepwise =} \ConstantTok{TRUE}\NormalTok{)}
\NormalTok{  fc\_ho    }\OtherTok{\textless{}{-}} \FunctionTok{forecast}\NormalTok{(m\_ho, }\AttributeTok{h =}\NormalTok{ holdout\_n)}
  \FunctionTok{mean}\NormalTok{(}\FunctionTok{abs}\NormalTok{((ts\_test }\SpecialCharTok{{-}}\NormalTok{ fc\_ho}\SpecialCharTok{$}\NormalTok{mean) }\SpecialCharTok{/}\NormalTok{ ts\_test)) }\SpecialCharTok{*} \DecValTok{100}
\NormalTok{\}, }\AttributeTok{error =} \ControlFlowTok{function}\NormalTok{(e) }\ConstantTok{NA\_real\_}\NormalTok{)}
\end{Highlighting}
\end{Shaded}

\subsection{5. Tasks Mode (Fixed Batch)}\label{tasks-mode-fixed-batch}

Tasks mode dispatches all chunks upfront using a Snowflake Task DAG.
Each task launches an SPCS job service that processes one chunk of
\texttt{unit\_ids}.

\begin{Shaded}
\begin{Highlighting}[]
\NormalTok{tasks\_out }\OtherTok{\textless{}{-}} \FunctionTok{parallel\_lab\_run\_tasks\_demo}\NormalTok{(conn, cfg,}
  \AttributeTok{run            =} \ConstantTok{TRUE}\NormalTok{,}
  \AttributeTok{arima\_stepwise =} \ConstantTok{TRUE}\NormalTok{,}
  \AttributeTok{forecast\_h     =} \DecValTok{30}\DataTypeTok{L}\NormalTok{,}
  \AttributeTok{save\_models    =} \ConstantTok{TRUE}
\NormalTok{)}
\end{Highlighting}
\end{Shaded}

\subsubsection{How it works}\label{how-it-works}

\begin{enumerate}
\def\labelenumi{\arabic{enumi}.}
\tightlist
\item
  \texttt{registerDoSnowflake(conn,\ mode\ =\ "tasks",\ ...)} sets up
  the backend.
\item
  The \texttt{foreach} loop iterates over \texttt{unit\_ids};
  doSnowflake chunks them into \texttt{chunks\_per\_job} groups.
\item
  Each chunk becomes an SPCS job service that:

  \begin{itemize}
  \tightlist
  \item
    Bulk-reads its partition of \texttt{SERIES\_EVENTS} via the
    \texttt{data\_query} spec
  \item
    Runs \texttt{mclapply} over the units within the container
  \item
    Writes \texttt{.rds} model files to
    \texttt{@STAGE/models/\textless{}run\_id\textgreater{}/\textless{}unit\_id\textgreater{}.rds}
  \item
    Writes result \texttt{.rds} back to
    \texttt{@STAGE/\textless{}job\textgreater{}/results/result\_\textless{}chunk\textgreater{}.rds}
  \end{itemize}
\item
  The orchestrator polls until all chunks complete, then collects
  results.
\end{enumerate}

\subsubsection{Chunk sizing}\label{chunk-sizing}

Balance parallelism vs overhead:

\begin{itemize}
\tightlist
\item
  \textbf{Too few chunks} (e.g.~1): all work in one container, no
  parallelism across nodes.
\item
  \textbf{Too many chunks} (e.g.~2000, one per SKU): excessive
  job-launch overhead.
\item
  \textbf{Sweet spot}: 4--20 chunks for 2000 units. Each container gets
  100--500 SKUs and uses \texttt{mclapply} for intra-node parallelism.
\end{itemize}

\subsection{6. Queue Mode (Time-Boxed
Streaming)}\label{queue-mode-time-boxed-streaming}

Queue mode uses a Hybrid Table as a work queue. A \textbf{dynamic
feeder} inserts rows as workers pull them, up to a time budget.

\begin{Shaded}
\begin{Highlighting}[]
\NormalTok{queue\_out }\OtherTok{\textless{}{-}} \FunctionTok{parallel\_lab\_run\_queue\_timeboxed}\NormalTok{(conn, cfg,}
  \AttributeTok{run            =} \ConstantTok{TRUE}\NormalTok{,}
  \AttributeTok{time\_limit\_sec =}\NormalTok{ tasks\_out}\SpecialCharTok{$}\NormalTok{elapsed,   }\CommentTok{\# match Tasks wall time}
  \AttributeTok{model\_contest  =} \ConstantTok{TRUE}\NormalTok{,                }\CommentTok{\# ARIMA + ETS + TBATS contest}
  \AttributeTok{forecast\_h     =}\NormalTok{ tasks\_out}\SpecialCharTok{$}\NormalTok{params}\SpecialCharTok{$}\NormalTok{forecast\_h,}
  \AttributeTok{n\_chunks       =} \DecValTok{40}\DataTypeTok{L}\NormalTok{,}
  \AttributeTok{unit\_ids       =}\NormalTok{ worst\_skus           }\CommentTok{\# worst{-}first from Tasks}
\NormalTok{)}
\end{Highlighting}
\end{Shaded}

\subsubsection{Dynamic feeder pattern}\label{dynamic-feeder-pattern}

Unlike Tasks, the Queue does \textbf{not} dispatch all work upfront. The
feeder loop:

\begin{enumerate}
\def\labelenumi{\arabic{enumi}.}
\tightlist
\item
  Checks how many chunks are \texttt{PENDING} in the queue.
\item
  If slots are available and the time budget hasn't expired, inserts the
  next chunk.
\item
  Launches SPCS workers that claim rows via compare-and-swap (CAS):
  \texttt{UPDATE\ ...\ SET\ STATUS\ =\ \textquotesingle{}RUNNING\textquotesingle{}\ WHERE\ STATUS\ =\ \textquotesingle{}PENDING\textquotesingle{}\ LIMIT\ 1}.
\item
  When the feeder's time budget expires, it stops feeding but waits for
  in-flight chunks to finish.
\end{enumerate}

This gives a natural ``time-boxed'' behavior: you spend exactly as much
wall time as you budget, processing as many chunks as workers can
handle.

\subsubsection{Accuracy-ranked feeding}\label{accuracy-ranked-feeding}

By ranking SKUs worst-first (highest holdout MAPE from the Tasks run),
the time budget is spent where it has the greatest impact:

\begin{Shaded}
\begin{Highlighting}[]
\NormalTok{worst\_skus }\OtherTok{\textless{}{-}} \FunctionTok{parallel\_lab\_rank\_skus\_by\_accuracy}\NormalTok{(}
  \AttributeTok{run\_results =}\NormalTok{ tasks\_out,}
  \AttributeTok{metric      =} \StringTok{"mape"}\NormalTok{,}
  \AttributeTok{descending  =} \ConstantTok{TRUE}
\NormalTok{)}
\end{Highlighting}
\end{Shaded}

\subsubsection{Model contest}\label{model-contest}

When \texttt{model\_contest\ =\ TRUE}, each SKU's worker expression runs
a holdout- validated competition:

\begin{enumerate}
\def\labelenumi{\arabic{enumi}.}
\tightlist
\item
  Split the daily series into training (\textasciitilde1065 days) and
  test (30 days).
\item
  Fit three candidates:

  \begin{itemize}
  \tightlist
  \item
    \texttt{auto.arima(stepwise\ =\ FALSE)} (exhaustive ARIMA search)
  \item
    \texttt{ets()} (exponential smoothing)
  \item
    \texttt{tbats()} (Trigonometric seasonality, Box-Cox, ARMA errors,
    Trend, Seasonal --- handles multiple seasonal periods natively,
    e.g.~weekly + annual cycles in daily data)
  \end{itemize}
\item
  The candidate with the lowest test-set MAPE wins.
\item
  The winner is refit on the full series for the production model.
\end{enumerate}

TBATS is particularly well-suited to this contest because the daily
synthetic data has \textbf{two nested seasonal cycles} (weekly period =
7, annual period = 365.25) that ARIMA and ETS cannot capture
simultaneously.

\textbf{ADAPT}: replace the three candidate model families with whatever
makes sense for your domain (Prophet, XGBoost, LSTM, etc.).

\subsection{7. Metrics and Comparison}\label{metrics-and-comparison}

\subsubsection{Persisting metrics}\label{persisting-metrics}

After each run, save per-unit accuracy metrics to the
\texttt{TRAINING\_METRICS} table:

\begin{Shaded}
\begin{Highlighting}[]
\FunctionTok{parallel\_lab\_save\_metrics}\NormalTok{(conn, cfg, tasks\_out)}
\FunctionTok{parallel\_lab\_save\_metrics}\NormalTok{(conn, cfg, queue\_out)}
\end{Highlighting}
\end{Shaded}

This writes
\texttt{(MODEL\_RUN\_ID,\ UNIT\_ID,\ ARIMA\_MODEL,\ AICC,\ RMSE,\ MAPE,\ MASE,\ N\_OBS)}
for every successfully fitted unit. The \texttt{MAPE} column stores the
\textbf{holdout} (out-of-sample) MAPE --- not the training MAPE.

\subsubsection{Why holdout MAPE?}\label{why-holdout-mape}

Training MAPE reflects how well the model fits historical data. Holdout
MAPE reflects how well it \textbf{predicts unseen data}. When comparing
a simple stepwise ARIMA (Tasks) against a model contest (Queue),
training MAPE will look similar because both methods fit the full
series. The difference only shows up in out-of-sample performance.

\subsubsection{Plotting}\label{plotting}

Single-run histogram:

\begin{Shaded}
\begin{Highlighting}[]
\FunctionTok{parallel\_lab\_plot\_accuracy}\NormalTok{(tasks\_out, }\AttributeTok{metric =} \StringTok{"mape"}\NormalTok{)}
\end{Highlighting}
\end{Shaded}

Faceted comparison (Tasks vs Queue, matched SKUs only):

\begin{Shaded}
\begin{Highlighting}[]
\FunctionTok{parallel\_lab\_plot\_accuracy\_comparison}\NormalTok{(}
\NormalTok{  tasks\_out, queue\_out,}
  \AttributeTok{metric =} \StringTok{"mape"}\NormalTok{,}
  \AttributeTok{matched\_only =} \ConstantTok{TRUE}
\NormalTok{)}
\end{Highlighting}
\end{Shaded}

The \texttt{matched\_only\ =\ TRUE} parameter restricts the comparison
to SKUs that appear in both runs, making the distributions directly
comparable.

\subsection{8. Model Registry}\label{model-registry}

\subsubsection{Registering N models as one
version}\label{registering-n-models-as-one-version}

After training, each SKU's \texttt{.rds} model file lives at
\texttt{@STAGE/models/\textless{}run\_id\textgreater{}/\textless{}unit\_id\textgreater{}.rds}.
The stage's \texttt{DIRECTORY} table tracks all files. The
\texttt{MODEL\_INDEX} view extracts \texttt{RUN\_ID} and
\texttt{PARTITION\_KEY} from the file paths.

\texttt{sfr\_log\_many\_model()} registers a \textbf{CustomModel
aggregator} in Model Registry. The aggregator:

\begin{itemize}
\tightlist
\item
  Carries a \texttt{model\_index.json} mapping \texttt{UNIT\_ID} → stage
  path
\item
  Exposes \texttt{@partitioned\_api\ predict()} for \texttt{run\_batch}
  (TABLE\_FUNCTION)
\item
  Exposes \texttt{@inference\_api\ predict\_single()} for SPCS service
  (FUNCTION)
\item
  Registers in O(1) time (\textasciitilde8--13s regardless of partition
  count)
\end{itemize}

\begin{Shaded}
\begin{Highlighting}[]
\NormalTok{reg\_tasks }\OtherTok{\textless{}{-}} \FunctionTok{parallel\_lab\_register\_models}\NormalTok{(conn, cfg,}
  \AttributeTok{model\_run\_id =}\NormalTok{ tasks\_out}\SpecialCharTok{$}\NormalTok{model\_run\_id,}
  \AttributeTok{forecast\_h   =}\NormalTok{ tasks\_out}\SpecialCharTok{$}\NormalTok{params}\SpecialCharTok{$}\NormalTok{forecast\_h}
\NormalTok{)}
\end{Highlighting}
\end{Shaded}

\subsubsection{The Ray shuffle bug
(SNOW-2193288)}\label{the-ray-shuffle-bug-snow-2193288}

Model Registry's \texttt{run\_batch} uses Ray internally for partitioned
dispatch. When the number of unique partition keys is much less than
\textasciitilde200, Ray's shuffle mechanism can fail. \textbf{Always
ensure \textgreater= 200 partition keys} when using \texttt{run\_batch}.
The reference implementation warns if this threshold is not met.

\subsection{9. Batch Inference}\label{batch-inference}

\subsubsection{Server-side partitioned
inference}\label{server-side-partitioned-inference}

\texttt{sfr\_run\_batch()} launches an SPCS batch job that:

\begin{enumerate}
\def\labelenumi{\arabic{enumi}.}
\tightlist
\item
  Partitions the input DataFrame by \texttt{UNIT\_ID}.
\item
  For each partition, the \texttt{RManyModelAggregator} downloads the
  \texttt{.rds} model from stage, loads it via \texttt{rpy2}, and runs
  \texttt{forecast::forecast(model,\ h=...)}.
\item
  Returns a DataFrame of
  \texttt{(PARTITION\_KEY,\ PERIOD,\ FORECAST,\ LOWER\_80,\ \ \ \ UPPER\_80,\ LOWER\_95,\ UPPER\_95,\ STATUS)}.
\end{enumerate}

\begin{Shaded}
\begin{Highlighting}[]
\NormalTok{fc\_df }\OtherTok{\textless{}{-}} \FunctionTok{parallel\_lab\_run\_inference}\NormalTok{(conn, cfg,}
  \AttributeTok{model\_run\_id =}\NormalTok{ tasks\_out}\SpecialCharTok{$}\NormalTok{model\_run\_id,}
  \AttributeTok{forecast\_h   =}\NormalTok{ tasks\_out}\SpecialCharTok{$}\NormalTok{params}\SpecialCharTok{$}\NormalTok{forecast\_h}
\NormalTok{)}
\end{Highlighting}
\end{Shaded}

The results are written to the \texttt{FORECAST\_RESULTS} table with
columns:
\texttt{RUN\_ID,\ UNIT\_ID,\ HORIZON,\ FORECAST\_DATE,\ POINT\_FORECAST,\ LO\_80,\ HI\_80,\ LO\_95,\ HI\_95,\ STATUS}.

\subsubsection{Input contract}\label{input-contract}

The input DataFrame needs at minimum a \texttt{UNIT\_ID} column (the
partition key). Additional columns are ignored by the aggregator --- it
loads the model from stage rather than re-training.

\subsection{10. Monitoring}\label{monitoring}

The reference implementation includes two Streamlit dashboards:

\begin{itemize}
\tightlist
\item
  \textbf{\texttt{streamlit\_benchmark\_monitor.py}}: live progress for
  Tasks vs Queue benchmark runs, filtered by \texttt{JOB\_ID}.
\item
  \textbf{\texttt{streamlit\_parallel\_demo\_monitor.py}}:
  general-purpose SPCS service and queue monitoring.
\end{itemize}

Deploy to Snowflake (Streamlit-in-Snowflake) with:

\begin{Shaded}
\begin{Highlighting}[]
\ExtensionTok{python}\NormalTok{ deploy\_streamlit\_monitor.py }\DataTypeTok{\textbackslash{}}
  \AttributeTok{{-}{-}connection{-}name}\NormalTok{ my\_demo\_jwt }\DataTypeTok{\textbackslash{}}
  \AttributeTok{{-}{-}private{-}key{-}file}\NormalTok{ \textasciitilde{}/.snowflake/keys/rsa\_key.p8}
\end{Highlighting}
\end{Shaded}

\textbf{Tip}: scope dashboard queries by \texttt{JOB\_ID} to avoid
showing cumulative counts from previous runs.

\subsection{11. Lessons Learned}\label{lessons-learned}

\subsubsection{JWT expiry and stale
sessions}\label{jwt-expiry-and-stale-sessions}

\texttt{sfr\_connect()} auto-detects Workspace Notebook sessions. When
running in RStudio with explicit \texttt{name} / \texttt{account}
parameters, it skips auto-detection to avoid wrapping a stale, expired
Python session. If you see \texttt{JWT\ token\ is\ invalid} after a
timeout, explicitly reconnect:

\begin{Shaded}
\begin{Highlighting}[]
\NormalTok{conn }\OtherTok{\textless{}{-}}\NormalTok{ snowflakeR}\SpecialCharTok{::}\FunctionTok{sfr\_connect}\NormalTok{(}
  \AttributeTok{name =} \StringTok{"my\_demo\_jwt"}\NormalTok{,}
  \AttributeTok{private\_key\_file =} \StringTok{"\textasciitilde{}/.snowflake/keys/rsa\_key.p8"}
\NormalTok{)}
\NormalTok{dbi\_con }\OtherTok{\textless{}{-}}\NormalTok{ snowflakeR}\SpecialCharTok{::}\FunctionTok{sfr\_dbi\_connection}\NormalTok{(conn)}
\end{Highlighting}
\end{Shaded}

\subsubsection{RStudio Connections Pane}\label{rstudio-connections-pane}

Call \texttt{sfr\_dbi\_connection(conn)} after \texttt{sfr\_connect()}
to populate the Connections Pane with database catalog info. If schemas
or tables don't appear, disconnect in the pane and reconnect --- RStudio
caches connection observer closures.

\subsubsection{Setting the right ts
frequency}\label{setting-the-right-ts-frequency}

Use \texttt{ts(data,\ frequency\ =\ 7)} for daily data with a weekly
cycle, or \texttt{frequency\ =\ 12} for monthly data. Without the
correct frequency, \texttt{auto.arima} cannot learn seasonal patterns
and forecasts will be flat. For daily data with both weekly and annual
cycles, TBATS auto-detects multiple seasonal periods from the
\texttt{frequency} attribute.

\subsubsection{Chunk count parity}\label{chunk-count-parity}

When comparing Tasks vs Queue, use chunk counts that give comparable
total work. The Queue's dynamic feeder means fewer chunks may complete
within the time budget --- that's by design (time-boxed, not
work-boxed).

\subsubsection{\texorpdfstring{\texttt{bquote} vs \texttt{quote} for
parameter
substitution}{bquote vs quote for parameter substitution}}\label{bquote-vs-quote-for-parameter-substitution}

The worker expression uses \texttt{bquote()} with \texttt{.()} splicing
to inject parameter values (e.g.~\texttt{.(forecast\_h)},
\texttt{.(arima\_stepwise)}) at definition time. Using \texttt{quote()}
would capture the variable \emph{names}, not their \emph{values},
causing failures when the expression runs in a remote container where
those variables don't exist.

\begin{Shaded}
\begin{Highlighting}[]
\NormalTok{forecast\_expr }\OtherTok{\textless{}{-}} \FunctionTok{bquote}\NormalTok{(\{}
\NormalTok{  model }\OtherTok{\textless{}{-}} \FunctionTok{auto.arima}\NormalTok{(ts\_full, }\AttributeTok{stepwise =}\NormalTok{ .(arima\_stepwise))}
\NormalTok{  fc    }\OtherTok{\textless{}{-}} \FunctionTok{forecast}\NormalTok{(model, }\AttributeTok{h =}\NormalTok{ .(forecast\_h))}
  \CommentTok{\# ...}
\NormalTok{\})}
\end{Highlighting}
\end{Shaded}

\subsection{Further Reading}\label{further-reading}

\begin{itemize}
\tightlist
\item
  \textbf{\texttt{vignette("parallel-dosnowflake")}}: the low-level
  \texttt{registerDoSnowflake} API reference (foreach backend, generic
  SPCS executor, crew/mirai workers).
\item
  \textbf{Reference implementation}:
  \texttt{internal/doSnowflake/demos/parallel\_forecast/} contains the
  complete runnable demo with configs, deploy scripts, and tests.
\item
  \textbf{Function reference}: \texttt{?sfr\_log\_many\_model},
  \texttt{?sfr\_run\_batch}, \texttt{?sfr\_model\_registry},
  \texttt{?registerDoSnowflake}.
\end{itemize}

\end{document}
